\subsection{Uma reserva de sete unidades}
\label{subsec:reserva-sete-unidades}

Alice tem, na mesma conta, três saídas confirmadas e elegíveis de 12,
8 e 3 unidades. Conhece as chaves necessárias e nenhuma das três
está gasta ou congelada. Pedem-lhe que demonstre uma reserva mínima de
7. A conta inteira vale \(12+8+3=23\); uma saída de 8 basta para a
declaração. Também bastaria a de 12. Convém seguir a escolha da carteira,
pois o verificador vai conhecer o valor inteiro de cada saída que Alice
incluir.

Ordenadas as saídas, ficam 12, 8 e 3. Enquanto a segunda da lista bastar
por si, a carteira pode retirar a primeira e conservar a possibilidade de
responder com uma só saída. Neste caso retira 12, porque \(8\geq7\).
Observa em seguida que \(3<7\) e fica com a primeira que resta, 8.
O verificador examina a prova relativa a essa saída, consulta a imagem
de chave e obtém \(\texttt{total}=8\), \(\texttt{spent}=0\). A condição
que lhe interessa é \(8-0\geq7\).

A diferença entre atingir e exceder o mínimo vê-se com um novo pedido
de 8. A saída escolhida continua a ser 8 e a desigualdade verifica-se
agora com igualdade. Para um pedido de 13, em contrapartida, falta
sempre alguma coisa quando se mostra uma única saída. A carteira começa
da maior, 12, junta-lhe 8 e pára com 20. Reunamos os casos, enquanto o
estado das saídas permanece o mesmo.

\begin{center}
\begin{tabular}{c|c|c|c}
mínimo pedido & saídas incluídas & total revelado & excedente\\
\hline
3  & 3     &  3 & 0\\
7  & 8     &  8 & 1\\
8  & 8     &  8 & 0\\
13 & 12, 8 & 20 & 7\\
21 & 12, 8, 3 & 23 & 2
\end{tabular}
\end{center}

Para o mínimo de 13 haveria outro par possível, 12 e 3, de soma 15.
Ambos os pares expõem duas saídas. A regra aplicada escolhe o primeiro
par na ordem decrescente quando já sabe que nenhuma saída serve sozinha.
É fácil justificar por que obtém o menor número de saídas. Se
\(a_1\geq a_2\geq\cdots\geq a_n\), a maior soma que se pode formar
com \(k\) saídas é \(a_1+\cdots+a_k\). Enquanto essa soma for
inferior ao pedido, todo o conjunto de \(k\) saídas será insuficiente.
Contudo, uma vez atingido o número de saídas necessário, o excesso
revelado depende das parcelas escolhidas. As sete unidades de excesso
da quarta linha são uma consequência visível desse critério de selecção
\cite{monero-source-v16}.

\paragraph{A mesma prova vista em dois momentos.}
Guardemos a prova do mínimo de 13. Contém os elementos relativos às
saídas de 12 e 8, incluindo as duas imagens de chave. No primeiro
momento ambas figuram como disponíveis. A reserva observada para
esse conjunto é \(20-0=20\), suficiente para o pedido.

Suponhamos agora que Alice gastou a saída de 8, que a transacção de
despesa ficou confirmada e que os dois observadores consultam o mesmo
ramo já sincronizado. As assinaturas da prova que guardámos ainda
conferem. As saídas de criação ainda têm os mesmos montantes e as suas
assinaturas referem as mesmas chaves. A resposta à consulta da segunda
imagem mudou. Obtemos agora
\[
   \texttt{total}=20,\qquad \texttt{spent}=8,\qquad 20-8=12<13 .
\]
O documento recebido por Bob conservou cada byte. A condição de reserva
dependia também da informação de gasto disponível na altura da consulta.

A terceira saída, de 3, continuou fora daquela prova. Se ainda estiver
disponível, Alice dispõe, entre as saídas antigas, de \(12+3=15\).
Poderá gerar outra prova para o mesmo pedido. Bob, porém, ao examinar a
primeira, só tem justificação para contar os 12 que nela permanecem
disponíveis. A diferença entre 12 e 15 mede aqui precisamente o efeito
de seleccionar as saídas a revelar. Qualquer troco recebido na despesa
seria ainda uma saída nova, com a sua própria chave e a sua própria
imagem. Também essa saída pediria uma prova que a incluísse.

Do lado de quem recebe a prova, há assim uma pequena conta que vale a
pena conservar juntamente com a mensagem de contexto. Registam-se o
mínimo pedido, o conjunto que deu o total, o momento da consulta,
a ponta do ramo e o montante que o nó classificou como gasto. No
exemplo, os gastos considerados estavam já confirmados. A consulta de
uma carteira pode também contar como gasta uma imagem que o nó vê
numa transacção ainda em memória. Nesse caso pode mudar a resposta
mesmo antes de mudar a altura do ramo. O momento da consulta situa,
portanto, uma parte da informação que as assinaturas da reserva deixam
em aberto.
